\section{Preliminary}


\subsection{Incomplete Knowledge Graph Question Answering}
IKGQA is a generalization of the standard KGQA task. In KGQA, the goal is to predict the correct answer entities $A_q \subseteq E$ for a given natural language question $q$, based on a complete knowledge graph $G = (E, R, T)$. Here, $E$ denotes the set of entities, $R$ denotes the set of relations, and $T = { (e_h, r, e_t) \mid e_h, e_t \in E, r \in R }$ represents the set of factual triples. Each triple consists of a head entity $e_h$, a relation $r$, and a tail entity $e_t$. This task assumes that all topic entities $T_q \subseteq E$ and the necessary relation paths are present in $G$. Formally, KGQA can be defined as a function ${f: (q, G) \mapsto A_q}$.

In real-world applications, however, KGs are often incomplete due to limitations in construction and maintenance. To address this, IKGQA relaxes the requirement for full KG coverage by enabling the reasoning process to leverage external knowledge—either retrieved from textual sources or derived from the internal knowledge embedded in LLMs. Following GoG~\cite{xu2024generate}, IKGQA can be formulated as ${f: (q, G, \mathcal{R}) \mapsto A_q}$, where $\mathcal{R}$ denotes external knowledge that supplements the incomplete KG $G$ during the inference process.


\subsection{Notation for LLM Modules}  
We denote the outputs of task-specific LLM calls using the notation $\text{LLM}_{\text{role}}(\cdot)$, where the subscript indicates the role or function performed by the LLM in the inference process. For example, $\text{LLM}_{\text{plan}}$ represents the planning module for sub-question decomposition.



